\chapter{Basic Use}
\label{chap:basicUse}

\begin{chapabstract}
In this chapter we present a minimum working example of a fully-formatted thesis chapter. Most of the core fundamentals will be introduced. By the end of this chapter you should be well on your way to making your own thesis.
\end{chapabstract}

Writing text in \LaTeX is quite simple. By using the appropriate commands, we can make the font style \textbf{bold}, \textit{italic}, \texttt{typewritten}, \textsc{small cap}, `things in quotes', ``things in double quotes'', hyphenate-things and so forth. We can cite previous studies, such as \cite{testington2011effects}, or any others you have saved in your bibtex file \citep[\eg][]{neddis2009best, jones1986coping}. As you will see in \sref{sec:exampleSection}, we can also reference other parts of this document.


% ==============================================================================================
\section{Graphics, tables and equations}
\label{sec:exampleSection} % create a reference handle, for reference elsewhere

%------------------------------------------------------------------------------
\subsection{A figure}
An example figure is given In \fref{fig:exampleFigure}.

\begin{figure}[H]
\includegraphics[scale=.3]{figs/exampleFigure}
\caption[short caption to appear in list of figures]
{
long title to accompany the figure. Note that since we didn't specify the file extension, the class will automatically search for an appropriate file with the stem `exampleFigure'. If you are using the PdfLatex engine, it will search for .pdf or bitmap files (\eg~.png, .jpg). If you are using the Latex engine, it will search for vector files (\eg~.eps, .svg). The recommended approach is to use .pdf figures.
}
\label{fig:exampleFigure}
\end{figure}


%------------------------------------------------------------------------------
\subsection{A table}
An example table is given In \fref{tab:exampleTable}.

\begin{table}[H]
\centering
\begin{tabular}{lccccr}
\toprule
id&score&q-factor&y-factor&fraud-rate&notes\\
\midrule
PJ	& 70.4 & 9	& 9 	& 22.1 & \textit{highly questionable}\\
NJ	& 46.9 & 8	& 11 & 17.9 &\\
CT& 90.0 & 13	& 6 	& 19.2 & \textit{exemplary}\\
\bottomrule
\end{tabular}
\caption[short caption to appear in list of tables]
{
long title to accompany the table.
}
\label{tab:exampleTable}
\end{table}

%------------------------------------------------------------------------------
\subsection{An equation}
Numbered equations can be entered into Matlab thus:
\begin{equation}
c^2 = a^2 + b^2 \textrm{.}
\label{eq:exampleEquation}
\end{equation}
The equation environment was used to create \eref{eq:exampleEquation}. We can also type maths inline, within the body of the text. This can be of useful, for example, when reporting statistics [$F_{(2,3)} = 4.2, p = .002, \eta^{2}_{p} = 0.7$].


% ==============================================================================================
\section{Conclusions}

By editing this chapter you should now be able to:
\begin{enumerate}[(1)]
\item Insert text.
\item Divide a chapter into separate sections
\item Insert graphics, tables and equations
\item Link to sections and graphics
\item Add key components, such as abstracts and acknowledgments, appendices, etc.
\end{enumerate}


% ==============================================================================================
\begin{chapacknowledgements}
This work was supported by the Medical Research Council, UK (Grant: U135097130).
\end{chapacknowledgements}

%% EOF